\documentclass[10pt,a4paper]{amsart}
\usepackage[margin=19mm]{geometry}
\usepackage{amsmath,amssymb,mathtools}
\usepackage[scr=rsfs,cal=euler]{mathalfa}
\usepackage{xfrac}
\usepackage[unicode,hidelinks,bookmarks=false]{hyperref}
\newcommand{\ie}{i.e.\ }
\newcommand{\cf}{cf.\ }
\newcommand{\resp}{respectively\ }
\newcommand{\Subsection}{\subsection}
\providecommand{\nobreakdash}{\nobreak\hbox{-}}
\setlength{\emergencystretch}{2em}
\multlinegap0pt
\usepackage{braket}
\usepackage{amssymb}
\usepackage{dsfont}
\usepackage{xcolor}
\newtheorem{proposition}{Proposition}
\title{Deterministic identification for Bernoulli channels and related channels with continuous input}
\author{Pau Colomer, Christian Deppe, Holger Boche, Andreas Winter}
\date{Source: arXiv:2605.05168v1 (CC BY 4.0)}
\begin{document}
\maketitle

\noindent\emph{Source and licence.} Deterministic identification for Bernoulli channels and related channels with continuous input. Version arXiv:2605.05168v1 (CC BY 4.0), distributed by arXiv under the Creative Commons Attribution 4.0 International licence, http://creativecommons.org/licenses/by/4.0/, as recorded in the arXiv licence metadata for this exact version and shown on the abstract page https://arxiv.org/abs/2605.05168v1. Full licence notice: \texttt{licenses/arxiv-2605.05168-CC-BY-4.0.txt}.\par\smallskip

\noindent\textit{Editorial scope.} Counted unit: the article's proposition proposition:concentration (source0.tex lines 471--476). The article's complete original proof of it is delivered with it (source0.tex lines 477--494, one \texttt{proof} environment of 18 lines). No mathematics is written, repaired or re-derived: the statement and the proof are the article's own lines, in the article's own notation, and no retained line is substituted or re-flowed.  Retained uncounted as the source-local context the proof needs: lines 422--427.  One declared adaptation: exactly 1 ancillary strings inside the retained spans are omitted (image-inclusion commands and, where the source repeats a label, the repeated label definitions) and no other line differs from the source; the figure environments, with their placement, captions and labels, are kept, and the package records the omitted strings in fidelity receipts bound to the affected bodies.  No retained line contains a citation, so the wrapper carries no bibliography and no bibliography entry.  The retained lines contain the article's own diagram code, which the wrapper typesets with the standard tikz packages; no image file is delivered and the package declares no asset dependency.  Editorial formatting only, in addition: an AMS \texttt{amsart} preamble that restates the article's own declarations and macros used by the retained lines, namely the environments \texttt{proposition} on separate counters, together with the standard packages those lines need.  The retained environments print in this document as Proposition 1 in order of appearance, whereas the article prints the same items with the numbering of its own full document.  The complete native source of the article as distributed in the arXiv e-print for this exact version is bundled unchanged as \texttt{sources/199929/source0.tex}.
\begin{figure}
    \centering
    
    \caption{Schematic structure of a $d$-angle-dense arrangement at layer $\ell$ around the point $\vec o_{s^{\ell-1}}$ (which, at the same time, is a point in the $\ell-1$-layer). Observe that the orthogonal projection of the $\ell$ layer point $\vec o_{\tilde s^\ell}$ onto the vector $\vec v_{s_{\ell-1}\to s_\ell}$ is at least $d$. The two red lines, respectively orthogonal to the vectors $\vec v_{s_{\ell-1}\to s_\ell}$ and $\vec v_{s_{\ell-1}\to \tilde s_\ell}$ represent the regions where the $\ell+1$ layer will be arranged. Notice that any point in the $\ell+1$ layer around $\vec o_{\tilde s^\ell}$ will have an orthogonal projection onto $\vec v_{s_{\ell-1}\to s_\ell}$ at a distance at most $\Delta$ from $\Pi_{s_{\ell-1}\to s_\ell}\vec o_{\tilde s^\ell}$ (see Proposition~\ref{proposition:concentration}).}
    \label{fig:subspaces}
\end{figure}

\begin{proposition}\label{proposition:concentration}
Let $\vec o_{s^L}\neq\vec o_{\tilde s^L}$ be two different last-layer elements (primitive code words) which differ at some layer $\ell$, that is, $s_\ell\in s^L$ is different from $\tilde s_\ell\in \tilde s^L$. Then,
\[
\|\Pi_{s_{\ell-1}\to s_\ell}\vec o_{\tilde s^L}-\Pi_{s_{\ell-1}\to s_\ell}\vec o_{s^L}\|\geq d-\Delta\geq d-\sqrt{2Ld}.
\]    
\end{proposition}

\begin{proof}
One just needs to notice that all the vectors $\vec v_{s_{\ell-1}\to s_\ell}$ corresponding to the same word, say $\vec o_{\tilde s^L}$ are mutually orthogonal. Hence, by Pythagoras, the distance between the code word $\vec o_{\tilde s^L}$ and its arrangement centre $\vec o_{\tilde s^\ell}$ at layer $\ell$ satisfies
\begin{equation}
\|\vec o_{\tilde s^L}-\vec o_{\tilde s^\ell}\|^2
\leq\sum_{j=\ell}^L \|\vec v_{s_{j-1}\to s_j}\|^2
=\sum_{j=\ell}^L r_j^2.
\end{equation}
Then, its projection onto the vector $\vec v_{s_{\ell-1}\to s_\ell}$ can be separated from the projection $\Pi_{s_{\ell-1}\to s_\ell}\vec o_{\tilde s^\ell}$ by at most a distance (cf.~Figure~\ref{fig:subspaces})
\begin{equation}\label{eq:delta_def}
    \Delta:=\sin\theta_\ell\sqrt{\sum_{j=\ell+1}^L r_j^2}.
\end{equation}

In the worst case scenario, the two code word projections we are comparing differ already at $\ell=1$, and angles are the smallest possible, that is, $\cos\theta_\ell=1-d/r_\ell$, and thus $\sin\theta_\ell=\sqrt{\frac{2d}{r_\ell}-\frac{d^2}{r_\ell^2}}$. Then,
\begin{equation}
  \Delta\leq\sqrt{L}r_2\sin\theta_1=\sqrt{L\left(2d-\frac{d^2}{r_1}\right)}\leq\sqrt{2Ld}.
\end{equation}
Where in the second line we used that $r_2=\sqrt{r_1}$. Then, the proposition claim follows trivially.
\end{proof}

\end{document}
